\documentclass{article}
\usepackage{amsmath,amssymb}
\usepackage{CJKutf8}
\usepackage[a4paper,margin=22mm]{geometry}

\setlength{\parindent}{0pt}
\setlength{\parskip}{0.5em}

\newcommand{\mondai}[1]{%
  \par\medskip
  \noindent{\large #1}\quad\ignorespaces
}

\newcommand{\sectionbox}[1]{%
  \par\bigskip
  \noindent\fbox{\large #1}\par\nobreak\smallskip
}

\newcommand{\defitem}[1]{%
  \item #1
}

\begin{document}
\begin{CJK}{UTF8}{min}

\begin{center}
{\Large 第4章　後半 演習問題}\\
\end{center}

\sectionbox{定義}

\begingroup\small

\begin{itemize}
\setlength{\itemsep}{0.15em}
\setlength{\parsep}{0pt}
\setlength{\topsep}{0.2em}
\setlength{\leftmargin}{1.5em}

\defitem{$\forall x\in X\,P(x)$ は $X$ の任意の元 $x$ について $P(x)$ が成り立つことを表す。$\exists x\in X\,P(x)$ は $P(x)$ を満たす $x\in X$ が存在することを表す。全称命題が偽であることを示すには反例を一つ与えればよい。}

\defitem{集合 $A,B$ に対して、$x\in A\cap B$ とは $x\in A$ かつ $x\in B$ であること、$x\in A\cup B$ とは $x\in A$ または $x\in B$ であることをいう。集合の等号 $A=B$ を示すには $A\subset B$ と $B\subset A$ を示せばよい。}

\defitem{写像 $f:X\to Y$ が単射とは $f(x)=f(x')$ ならば $x=x'$ が成り立つこと、全射とは任意の $y\in Y$ に対して $f(x)=y$ となる $x\in X$ が存在すること、全単射とは単射かつ全射であることをいう。$f:X\to Y$, $g:Y\to Z$ の合成は $g\circ f:X\to Z$, $(g\circ f)(x)=g(f(x))$ で定める。}

\defitem{空でない集合 $A\subset\mathbb{R}$ に対し、$u$ が上界とは任意の $a\in A$ について $a\leq u$ が成り立つことをいう。上界が存在するとき $A$ は上に有界であるという。上界のうち最小のものを上限といい、$\sup A$ と書く。}

\defitem{数列 $(a_n)$ が有界とは、集合 $\{a_n\mid n\geq 1\}$ が上にも下にも有界であることをいう。有界な数列に対し $A_n:=\{a_m\mid m\geq n\}$ とおくと、$\limsup_{n\to\infty}a_n:=\lim_{n\to\infty}\sup A_n$ を上極限という。}

\defitem{アルキメデスの公理とは、任意の $r\in\mathbb{R}$ に対して $n>r$ となる $n\in\mathbb{N}$ が存在することである。}

\defitem{数列 $(a_n)$ が Cauchy 列とは、任意の $\varepsilon>0$ に対して、ある $N$ が存在し、$m,n\geq N$ ならば $|a_m-a_n|<\varepsilon$ となることをいう。実数の完備性とは、実数内の任意の Cauchy 列が実数に収束することである。}

\defitem{実数の連続性の公理とは、空でない $A\subset\mathbb{R}$ が上に有界ならば、$\sup A$ が実数として存在することである。}


\end{itemize}

\endgroup

\sectionbox{問題}
\mondai{問1}次の命題が成り立つか判定せよ。理由も述べよ（証明でなくてよい）。
\begin{enumerate}
  \item $\exists x\in\mathbb{Z}\ \forall y\in\mathbb{Z},\ x\geq -y^2$
  \item $\forall x\in\mathbb{Z}\ \exists y\in\mathbb{Z},\ x=y^2$
\end{enumerate}

\mondai{問2}集合の規則に従って, $ A \cap (B \cup C) \subset (A \cap B) \cup (A \cap C)$を示せ。つまり
\[
x \in A \cap (B \cup C) \Rightarrow x \in (A \cap B) \cup (A \cap C)
\]
を示せ。ただし、命題のドモルガンの法則を用いてもよい。
\mondai{問3}(1) $f:X\to Y,g:Y\to Z$ が全単射であれば、合成 $g\circ f$ も全単射であることを証明せよ。
\par\noindent\hspace*{2.8em}(2) $f:X\to Y$ を写像とし、$A\subset X$ とする。このとき $A \subset f^{-1} (f(A))$ を示せ。また等号が成り立たない例を与えよ。

\mondai{問4}有界な数列 $(a_n)_{n=1}^{\infty}$ に対して
$A_n := \{a_m \mid m\geq n\}$ とおき、
$\sup_{m\geq n} a_m := \sup A_n$ とする。
この時、数列 $(a_n)_{n=1}^{\infty}$ の上極限
$\limsup_{n\to\infty} a_n := \lim_{n\to\infty}\sup_{m\geq n} a_m$
が存在することを示せ。なお、有界な単調減少列の極限が存在することを用いてよい。

\mondai{問5}命題 4.12 を証明せよ。すなわち、アルキメデスの公理と実数の完備性から、実数の連続性の公理を導け。
\end{CJK}\end{document}
